\section{Fuerza y masa}

La dinámica estudia la relación entre el movimiento de un cuerpo y las fuerzas que actúan sobre él.

\section{Primera ley de Newton (Inercia)}
\label{sec:primera-ley}
Todo cuerpo permanece en estado de reposo o de movimiento rectilíneo uniforme a menos que una fuerza externa neta actúe sobre él.

\section{Masa inercial y masa gravitacional}
\begin{itemize}
    \item \textbf{Masa inercial:} Mide la resistencia de un objeto a cambiar su estado de movimiento.
    \item \textbf{Masa gravitacional:} Mide la fuerza de atracción gravitacional entre dos cuerpos.
\end{itemize}

\section{Sistemas de referencia inerciales y no inerciales}
\begin{itemize}
    \item \textbf{Inerciales:} Se mueven con velocidad constante (no acelerados). En ellos se cumplen las leyes de Newton.
    \item \textbf{No inerciales:} Tienen aceleración. Se deben introducir fuerzas ficticias (como la fuerza centrífuga).
\end{itemize}

\section{Fuerzas comunes}
\begin{itemize}
    \item Peso: \(\vec{P} = m \cdot \vec{g}\)
    \item Normal: \(\vec{N}\), perpendicular a la superficie de contacto.
    \item Tensión: \(\vec{T}\), a lo largo de una cuerda o cable.
    \item Fricción: \(f = \mu \cdot N\) (estática o cinética).
\end{itemize}